\documentclass[12pt]{article}
\usepackage[utf8]{inputenc}
\usepackage[T1]{fontenc}
\usepackage{hyperref}
\title{Naluo Language Resources for K-12 Curriculum Development}
\author{Compiled by Claude Sonnet 5.5 (claude-sonnet-5-5) for Andr\'as Kornai}
\date{}
\begin{document}
\maketitle

\section{Introduction}
Naluo (ISO 639-3: ylo), also written Naruo, is a Loloish Sino-Tibetan language of the Yi people in Yunnan, China. Bradley (2004) puts the population at about 15,000, mostly in eastern Yongsheng County and southern Huaping County. Other names include Laluo, Shui Yi, Shuitian and Gan Yi. Ethnic-group counts differ (Asia Harvest about 37,000; Joshua Project 49,000) and describe the group rather than the language. Naluo should not be confused with Aluo (yna). No school materials were found.

\section{Primary Educational Resources}
None found.

\section{Linguistic and Reference Materials}
\begin{itemize}
\item Wikipedia, ``Naruo language''. \url{https://en.wikipedia.org/wiki/Naruo_language}
\item Wikipedia, ``Aluo language'', which is a different language and is listed only to avoid confusion. \url{https://en.wikipedia.org/wiki/Aluo_language}
\item Endangered Languages Project entry. \url{https://endangeredlanguages.com/lang/5132}
\end{itemize}

\section{Religious and Community Resources}
\begin{itemize}
\item Joshua Project, Naluo in China. \url{https://joshuaproject.net/people_groups/18620}
\item Asia Harvest people-group profile of the Naluo (PDF). \url{https://people-groups.asiaharvest.org/people-group-profiles/yunnan/naluo.pdf}
\end{itemize}

\section{Digital and Audio Resources}
None found.

\section{Curriculum Development Context}
The language is endangered in Yunnan and moribund or extinct in Sichuan. No orthography or school materials were found. Curriculum work would require new documentation and community partnership.

\end{document}